% Insertable IEEE manuscript subsection; not a complete or compiled paper.
% Merge references.bib and include all_weightings_table.tex where appropriate.
\subsection{Evaluation-distribution sensitivity}
We performed an exploratory sensitivity analysis of the saved validation
predictions from five stain-augmentation arms and three matched training seeds.
All models were evaluated on the same ordered 1,251 patches from seven patients.
The validation cohort had already informed early stopping and was selected using
annotations; these analyses are not independent test evaluations.

For target class masses $q_c$, patch-uniform weights were
$w_i=q_{y_i}/n_{y_i}$. For equal patient weighting within each class, we used
$w_i=q_{y_i}/(P_{y_i}n_{p_i,y_i})$, where $P_c$ is the number of represented
patients in class $c$ and $n_{p,c}$ counts their class-specific patches.
Global equal-patient weighting was analyzed separately. Ten class-mass scenarios
were derived from the training-only construction audit, alongside the observed
validation distribution and explicit uniform-class/patient scenarios, yielding
24 specified weightings. None was selected as a replacement primary endpoint.
No images were decoded or classifiers refitted for this analysis.

For weights summing to one, the weighted confusion matrix was
\begin{equation}
 C_{ab}=\sum_i w_i\mathbf{1}(y_i=a,\hat y_i=b),\qquad
 F_{1,\mathrm{macro}}=\frac{1}{5}\sum_{c=1}^5
 \frac{2C_{cc}}{\sum_b C_{cb}+\sum_a C_{ac}}.
\end{equation}
This is the mean of class-wise F1 scores, rather than the harmonic mean of
macro-precision and macro-recall~\cite{opitz2019macrof1}.
An independent sklearn implementation regenerated weights from per-record
denominators and class masses from the training-membership CSV, reproducing all
360 weighted confusion matrices and scalar scores within $2.8\times10^{-15}$.

Under observed validation frequencies, mean macro-F1 was 0.7131 for adaptive and
0.7350 for shuffled assignment, a difference of $-0.0219$.
Under equal class masses, the values were 0.7024 and 0.6920, respectively
(difference $+0.0104$). Fixed augmentation had the highest mean macro-F1 under
equal class masses. In contrast, equal global patient weighting retained the
original ordering and increased the adaptive-minus-shuffled deficit to
$-0.0575$. Every predefined scenario is reported in the accompanying table.
These correlated scenarios are not independent cohorts or hypothesis tests.

The observed-to-equal-class contrast shift was 0.03230 macro-F1, of which the
lymphoid class contributed 0.02243 (approximately 69.4\%). Both adaptive and
shuffled models recalled all seven lymphoid validation patches in each seed.
Their mean lymphoid precision nevertheless differed: 0.7583 versus 0.9583
under observed frequencies, and 0.9870 versus 0.9979 under equal class masses.
The changing false-positive mass explains the different F1 penalty; no additional
lymphoid cases were classified correctly by reweighting. Pure class-mass changes
preserved every class recall and balanced accuracy numerically.

Prevalence sensitivity and ranking changes are established properties of
classification evaluation~\cite{brabec2020imbalance,opitz2024metrics}.
Our result is a gastric, stain-augmentation-specific illustration and class-wise
decomposition, not a new metric or evidence of adaptive-controller superiority.
Equal class masses upweight the seven lymphoid cases from approximately 0.56\%
to 20\% of total evaluation mass; this is a hypothetical distribution, not an
estimate of clinical prevalence. The small, clustered, selected cohort and
validation-based model selection limit generalization. Reweighting does not
identify the effect of retraining with different purity, artifact or patient-cap
policies, which remains untested.
